\documentclass[10pt,a4paper]{amsart}
\usepackage[margin=19mm]{geometry}
\usepackage{amsmath,amssymb,mathtools}
\usepackage[scr=rsfs,cal=euler]{mathalfa}
\usepackage{xfrac}
\usepackage[unicode,hidelinks,bookmarks=false]{hyperref}
\newcommand{\ie}{i.e.\ }
\newcommand{\cf}{cf.\ }
\newcommand{\resp}{respectively\ }
\newcommand{\Subsection}{\subsection}
\providecommand{\nobreakdash}{\nobreak\hbox{-}}
\setlength{\emergencystretch}{2em}
\multlinegap0pt
\usepackage{amssymb}
\usepackage{mathtools}
\usepackage{bm}
\usepackage[shortlabels]{enumitem}
\newtheorem{lemma}{Lemma}
\usepackage{cleveref}
\crefname{lemma}{Lemma}{Lemmas}
\newcommand{\PP}{\mathbb{P}}
\title{Pareto-type finite-block optimality for source codes:\ a constrained Markov example}
\author{Stefano Della Fiore}
\date{Source: arXiv:2605.03552v1 (CC BY 4.0)}
\begin{document}
\maketitle

\noindent\emph{Source and licence.} Pareto-type finite-block optimality for source codes:\ a constrained Markov example. Version arXiv:2605.03552v1 (CC BY 4.0), distributed by arXiv under the Creative Commons Attribution 4.0 International licence, http://creativecommons.org/licenses/by/4.0/, as recorded in the arXiv licence metadata for this exact version and shown on the abstract page https://arxiv.org/abs/2605.03552v1. Full licence notice: \texttt{licenses/arxiv-2605.03552-CC-BY-4.0.txt}.\par\smallskip

\noindent\textit{Editorial scope.} Counted unit: the article's lemma lem:tail (source0.tex lines 419--428). The article's complete original proof of it is delivered with it (source0.tex lines 430--519, one \texttt{proof} environment of 90 lines). No mathematics is written, repaired or re-derived: the statement and the proof are the article's own lines, in the article's own notation, and no retained line is substituted or re-flowed.  Retained uncounted as the source-local context the proof needs: lines 394--409.  Nothing was omitted from any retained span, so the package carries no omission record.  No retained line contains a citation, so the wrapper carries no bibliography and no bibliography entry.  No retained line contains a figure environment, an image-inclusion command or drawing code, so no image is delivered and the package declares no asset dependency.  Editorial formatting only, in addition: an AMS \texttt{amsart} preamble that restates the article's own declarations and macros used by the retained lines, namely the environments \texttt{lemma} on separate counters, together with the standard packages those lines need.  The retained environments print in this document as Lemma 1 in order of appearance, whereas the article prints the same items with the numbering of its own full document.  The complete native source of the article as distributed in the arXiv e-print for this exact version is bundled unchanged as \texttt{sources/199795/source0.tex}.
\begin{lemma}[Conditional saving formula]\label{lem:cond}
Let $n\ge2$ and $x\in\{0,\dots,n-1\}$, and set $k:=n+1+x$. Define
\[
g_n(x):=\PP(I_n=1\mid X=x).
\]
Then
\begin{equation}\label{eq:gnV}
g_n(x)=\min\!\left\{1,\max\!\left\{0,\frac{V_n(k)}{N(k,n)}\right\}\right\}.
\end{equation}
Consequently,
\begin{equation}\label{eq:Pin-sum}
\PP(I_n=1)=\sum_{x=0}^{n-1}p_n(x)g_n(x),
\qquad
p_n(x)=\binom{n-1}{x}2^{-(n-1)}.
\end{equation}
\end{lemma}

\begin{lemma}[Tail saturation]\label{lem:tail}
Let $n\ge2$ and let $x\in\{0,\dots,n-1\}$. If
\[
x\ge \frac{n+1}{2},
\]
then
\[
g_n(x)=1.
\]
\end{lemma}

\begin{proof}
Fix $n\ge2$ and $x\ge (n+1)/2$, and set
\[
y:=n-1-x,
\qquad
k:=n+1+x,
\qquad
m:=k-2=n+x-1=x+2y.
\]
Using $N(k,\ell)=4U(k-2,k-\ell-1)$, we obtain
\[
N(k,n)=4U(m,x),
\qquad
W_{<n}(k)=4\sum_{j\ge x+1}U(m,j),
\qquad
S_k=4\sum_{j\ge0}U(m,j).
\]
Hence
\[
V_n(k)-N(k,n)
=2\left(\sum_{j=0}^{x-1}U(m,j)-\sum_{j\ge x+1}U(m,j)\right).
\]
Therefore it suffices to prove
\begin{equation}\label{eq:tail-target-v2}
\sum_{j=0}^{x-1}U(m,j)\ge \sum_{j\ge x+1}U(m,j).
\end{equation}

Since $m=x+2y$, the support of $j\mapsto U(m,j)$ is
\[
0\le j\le \Bigl\lfloor\frac m2\Bigr\rfloor=x+\Bigl\lfloor\frac y2\Bigr\rfloor,
\]
so the right-hand side of \eqref{eq:tail-target-v2} equals
\[
\sum_{r=1}^{\lfloor y/2\rfloor}U(m,x+r).
\]
We compare these terms with the left tail via
\[
R_r:=\frac{U(m,x-1-r)}{U(m,x+r)}
\qquad
\Bigl(0\le r\le \Bigl\lfloor\frac y2\Bigr\rfloor\Bigr).
\]
First,
\[
R_0=\frac{U(m,x-1)}{U(m,x)}=\frac{x(x+y+1)}{2(y+2)(y+1)}>1,
\]
because $x\ge y+2$.
Next, for $0\le r<\lfloor y/2\rfloor$, a direct computation gives
\[
\frac{R_{r+1}}{R_r}
=
\frac{(x+y-r)(x-r-1)(x+r+1)(x+y+r+2)}
{4(y-2r)(y-2r-1)(y+2r+3)(y+2r+4)}.
\]
The right-hand side is increasing in $x$, so it is bounded below by its value at the smallest admissible value $x=y+2$. Therefore
\[
\frac{R_{r+1}}{R_r}
\ge
\frac{(2y+2-r)(y-r+1)(y+r+3)(2y+r+4)}
{4(y-2r)(y-2r-1)(y+2r+3)(y+2r+4)}.
\]
Split this lower bound as
\[
\frac{(y-r+1)(y+r+3)}{(y-2r)(y+2r+4)}
\cdot
\frac{(2y+2-r)(2y+r+4)}{4(y-2r-1)(y+2r+3)}.
\]
Both factors are strictly larger than $1$, because
\[
(y-r+1)(y+r+3)-(y-2r)(y+2r+4)=3(r+1)^2>0,
\]
and
\[
(2y+2-r)(2y+r+4)-4(y-2r-1)(y+2r+3)=15r^2+30r+4y+20>0.
\]
Hence $R_{r+1}>R_r$, and therefore $R_r\ge R_0>1$ for every admissible $r$. Thus
\[
U(m,x-1-r)\ge U(m,x+r)
\qquad
\Bigl(0\le r\le \Bigl\lfloor\frac y2\Bigr\rfloor\Bigr).
\]
Summing over $r$ gives
\[
\sum_{j=0}^{x-1}U(m,j)
\ge
\sum_{r=0}^{\lfloor y/2\rfloor}U(m,x-1-r)
\ge
\sum_{r=1}^{\lfloor y/2\rfloor}U(m,x+r),
\]
which is exactly \eqref{eq:tail-target-v2}. Therefore $V_n(k)\ge N(k,n)$ and \Cref{lem:cond} yields $g_n(x)=1$.
\end{proof}

\end{document}
